% problem-000203 6 MAX相与MXene材料
\section*{6. MAX相与MXene材料}
MAX相是⼀⼤类具有层状结构的⾦属碳化物或氮化物的总称，其中M为Ti、V、Nb等前过渡⾦属，X为碳或氮，A为Al、Sn、Ge、Sb等�区元素。MAX的结构中，M原⼦形成理想的密置层，M层之间采取密堆积（可连续分布）与简单六⽅堆积（通常以单层呈现）按⼀定⽅式有序堆叠形成三维结构。结构中，X填充在M层密堆积形成的所有⼋⾯体空隙中，A则有序占据M层简单六⽅堆积所形成空隙的⼀半，将其中的A元素选择性除去，可以分离得到⼆维的层状结构，称为MXene。MXene层中，最外层M可进⼀步与⻧素、羟基等−1价端基T按1:1结合，形成端基T功能化的MXene (T-MXene)。因此，MXene的组成和结构多样且可调控，是当前⼆维材料的研究热点。

\subsection*{6-1}
若 MAX 相中 M 和 A 的原⼦数⽐为�，写出 MAX 相(O)、⼆维 MXene 层(P)和 T-MXene(Q)的组成通式。Q中T为−1价端基。

\subsection*{6-2}
某碳化物MAX相 $\mathrm { T i } _ { x } \mathrm { A l } _ { y } \mathrm { C } _ { z }$ 结构属六⽅晶系，Ti 层的排列⽅式为...ABCCBAABCCBA...。

\subsection*{6-2}
-1 写出 $\mathrm { T i } _ { x } \mathrm { A l } _ { y } \mathrm { C } _ { z }$ 晶胞的组成

\subsection*{6-2}
-2 已知 $\mathrm { T i } _ { x } \mathrm { A l } _ { y } \mathrm { C } _ { z }$ 晶胞参数 $a = 0 . 3 0 6 \ : \mathrm { n m }$ ，� =1.856 nm。将M层密堆积形成的⼋⾯体近似为正⼋⾯体，计算简单六⽅排布相邻层的间距�（单位：nm）。

6- $- 2 { \mathrm { - } } 3 \mathrm { T i } _ { x } { \mathrm { A l } } _ { y } { \mathrm { C } } _ { z }$ 结构中，碳原⼦处在Ti层密堆积形成的⼋⾯体中⼼。若将晶胞原点选在处于堆积中B层的Ti上，此时所有的Al恰好均处在�轴上。写出晶胞中所有碳原⼦的坐标参数。

\subsection*{6-2}
-4 将 $\mathrm { T i } _ { x } \mathrm { A l } _ { y } \mathrm { C } _ { z }$ 在HF溶液中超声处理，可选择去除其中的Al层，到⼆维结构的MXene⽚层 $\mathrm { T i } _ { x } \mathrm { C } _ { z }$ （反应1）， $\mathrm { T i } _ { x } \mathrm { C } _ { z }$ 进⼀步与F−反应形成F-MXene（反应2）。写出反应1和2的化学⽅程式。

\subsection*{6-2}
-5T-MXene的性质与端基T密切相关，利⽤熔融 $\mathrm { Z n C l _ { 2 } }$ $\mathrm { T i } _ { x } \mathrm { A l } _ { y } \mathrm { C } _ { z }$ 在 $5 5 0 ~ ^ { \circ } \mathrm { C }$ 反应5h，可以得到Cl为端基的 Cl-MXene（反应 3）。将 Cl-MXee 在 CsCl-KCl-LiCl 混合熔盐与 Li_{2}Se 反应 18 h，可以得到端基为Se的Se-MXene（反应4）。写出反应3和4的化学⽅程式。

\subsection*{6-2}
-6 对于6-2-4中⽤HF⽔溶液处理得到的F-MXene，能否发⽣与6-2-5中反应4类似的反应？简述原因。
